%%%%%%%%%%%%%%%%%%%%%%%%%%%%%%%%%%%%%%%%%%%%%%%%%%%%%%%%%%%%%%%%%%%%%%%%%%%%%%
%% Example Master's thesis using tvthesis class
%%
%% Compile with:
%%   pdflatex example-master
%%   bibtex example-master
%%   pdflatex example-master
%%   pdflatex example-master
%%
%% Class Options Summary (refer to the docs for other options):
%%  - english (default) / italian   : Selects thesis language for labels and headings
%%  - fem                           : (Italian only) Use feminine label ("Candidata")
%%  - lof                           : Automatically append List of Figures after TOC
%%  - lot                           : Automatically append List of Tables after TOC
%%%%%%%%%%%%%%%%%%%%%%%%%%%%%%%%%%%%%%%%%%%%%%%%%%%%%%%%%%%%%%%%%%%%%%%%%%%%%%
%%%%%%%%%%%%%%%%%%%%%%%%%%%%%%%%%%%%%%%%%%%%%%%%%%%%%%%%%%%%%%%%%%%%%%%%%%%%%%
\documentclass[master,italian,fem,lof,lot]{tvthesis}

% One of the following are allowed:
% Facoltà di Economia
% Dipartimento di Giurisprudenza
% Macroarea di Ingegneria
% Macroarea di Lettere e Filosofia
% Facoltà di Medicina e Chirurgia
% Macroarea di Scienze Matematiche, Fisiche e Naturali
\courseorganizer{Macroarea di Ingegneria}
\course{Laurea Magistrale in Ingegneria Informatica}
% The following is optional. Comment it out to remove from the title page.
\thesissubject{Sistemi Distribuiti}

\title{A Novel Approach to Distributed Computing}
\subtitle{Algorithms, Protocols, and Performance Analysis}
\author{Jane Doe}
\IDnumber{0123456}

\advisor{Prof.~Alice Smith}
\advisor{Prof.~Joe Bloggs}
\coadvisor{Prof.~Bob Johnson}
\coadvisor{Prof.~John Q. Public}
\coadvisor{Prof.~Robin Banks}

\AcademicYear{2025/2026}

\usepackage{lipsum}

\begin{document}

\maketitle

\frontmatter

\dedication{Alla mia famiglia}

\tableofcontents



\begin{acknowledgments}
Vorrei esprimere la mia sincera gratitudine al mio relatore per la guida preziosa e il continuo supporto durante tutto questo lavoro. Sono inoltre grato al mio co-relatore per i commenti illuminanti e l’incoraggiamento.

Un ringraziamento speciale alla mia famiglia e ai miei amici, il cui sostegno costante ha reso possibile questo percorso.
\end{acknowledgments}


\begin{abstract}
Questa tesi presenta un approccio innovativo al calcolo distribuito.
Proponiamo nuovi algoritmi e protocolli che migliorano le prestazioni
e la scalabilità rispetto alle soluzioni all’avanguardia. Un’ampia valutazione
sperimentale dimostra l’efficacia del nostro approccio in diversi scenari
di riferimento. I risultati mostrano miglioramenti significativi nella
capacità di elaborazione e nella latenza, mantenendo al contempo
solide garanzie di coerenza.
\end{abstract}


\mainmatter

\chapter{Introduzione}
\label{ch:introduction}

L'elaborazione distribuita è diventata un paradigma essenziale nell'ingegneria del software moderna. Con l'aumento della complessità e della scala dei sistemi, la necessità di algoritmi distribuiti efficienti diventa sempre più critica~\cite{tanenbaum2007}.

I contributi principali di questa tesi sono:
\begin{itemize}
    \item Un nuovo algoritmo di consenso.
    \item Un protocollo per la replicazione efficiente dei dati tra nodi distribuiti.
    \item Una valutazione sperimentale estensiva delle soluzioni proposte.
\end{itemize}

Il resto di questa tesi è organizzato come segue. Il Capitolo~\ref{ch:background} fornisce le nozioni di base necessarie. Il Capitolo~\ref{ch:proposal} presenta il nostro approccio proposto. Il Capitolo~\ref{ch:conclusions} trae le conclusioni e delinea il lavoro futuro.


\chapter{Concetti Fondamentali}
\label{ch:background}

In questo capitolo rivediamo i concetti fondamentali ed i lavori correlati che costituiscono la base della nostra ricerca.

\section{Sistemi Distribuiti}

Un sistema distribuito è un insieme di computer indipendenti che appare ai suoi utenti come un unico sistema coerente~\cite{tanenbaum2007}. Le principali sfide nella progettazione di tali sistemi includono:
\begin{enumerate}
    \item Raggiungere il consenso tra nodi distribuiti.
    \item Garantire tolleranza ai guasti e disponibilità.
    \item Mantenere la coerenza tra i dati replicati.
\end{enumerate}

\section{Algoritmi di Consenso}

Gli algoritmi di consenso sono componenti fondamentali per i sistemi distribuiti. Gli algoritmi più noti includono Paxos~\cite{lamport1998} e Raft~\cite{ongaro2014}.

\lipsum[1-2]


\chapter{Approccio Proposto}
\label{ch:proposal}

In questo capitolo presentiamo il nostro approccio proposto al calcolo distribuito
in ambienti eterogenei.

\section{Modello di Sistema}

Consideriamo un sistema distribuito composto da $n$ nodi interconnessi da
una rete. Ogni nodo $i$ mantiene uno stato locale $s_i$ e comunica
con gli altri nodi tramite passaggio di messaggi. Il sistema soddisfa le
seguenti proprietà:
%
\begin{equation}
\label{eq:consistency}
\forall i, j \in \{1, \ldots, n\}: si = sj \text{ in quiete}
\end{equation}

L’architettura ad alto livello del nostro sistema è mostrata in
Figura~\ref{fig:architecture}.

\begin{figure}[t]
    \centering
    \caption{Architettura ad alto livello del sistema distribuito proposto.}
    \label{fig:architecture}
    \includegraphics[width=0.7\textwidth]{example-image}
\end{figure}

La Tabella~\ref{tab:comparison} riassume le principali proprietà degli algoritmi
di consenso esistenti confrontate con il nostro approccio.

\begin{table}[t]
    \centering
    \begin{tabular}{lccc}
        \toprule
        \textbf{Algoritmo} & \textbf{Tolleranza ai guasti} & \textbf{Latenza} & \textbf{Throughput} \\
        \midrule
        Paxos      & $f < n/2$ & Alta   & Media \\
        Raft       & $f < n/2$ & Media  & Media \\
        Nostro Metodo & $f < n/2$ & Bassa & Alta  \\
        \bottomrule
    \end{tabular}
    \caption{Confronto tra algoritmi di consenso.}
    \label{tab:comparison}
\end{table}

\section{Progettazione dell’Algoritmo}

Il nostro algoritmo procede a turni. In ogni turno $r$, ogni nodo $i$:
\begin{enumerate}
    \item Trasmette in broadcast il proprio stato corrente $s_i^{(r)}$ a tutti gli altri nodi.
    \item Raccoglie gli stati da una maggioranza di nodi.
    \item Aggiorna il proprio stato in base alla funzione di fusione deterministica.
\end{enumerate}

\lipsum[3]

\chapter{Conclusioni}
\label{ch:conclusions}

In questa tesi, abbiamo presentato un approccio innovativo al calcolo distribuito.
I nostri risultati sperimentali dimostrano miglioramenti significativi rispetto alle soluzioni esistenti, sia in termini di throughput che di latenza.

Il lavoro futuro si concentrerà sull’estensione del nostro approccio per supportare implementazioni geo-distribuite e sull’indagine dell’applicazione di tecniche di apprendimento automatico per la selezione adattiva dei protocolli.



\bibliographystyle{plain}
\bibliography{references}

\end{document}
